\section{A Chow relation and its universal diagonal coefficient}
\label{sec:relation}

Our goal is the relation \(0=c[\Delta_X]+\Gamma\) of
\eqref{eq:intro-diagonal}, with \(\Gamma\) a sum of external products.
We first obtain a vanishing by equivariant degree, then decompose each
fixed contribution into a diagonal term and external products, retaining
both surface factors.

Throughout this Section Chow groups have rational coefficients.  We use
the equivariant perfect obstruction theory of Proposition~\ref{prop:quot-pot},
its fixed theory in Proposition~\ref{prop:fixed-pot}, and the fixed virtual
classes computed in Propositions~\ref{prop:fixed-pg0}
and~\ref{prop:fixed-auxiliary}.  Both settings have
$\chi(\mathcal O_X)=1$.  The torus parameter is initially the indeterminate
$t$, of equivariant codimension one; the weights on $L_1,L_2$ are $t,-t$.
Only after localization will we set $t=1/2$, obtaining
$w_1=1/2,w_2=-1/2$.

\subsection{The vanishing supplied by equivariant degree}

Fix $p\geq 0$ such that $v>p+2$, where
$v=2n+(A-K)D$ is the virtual dimension in~\eqref{eq:vdim}.
Let $Q$ also denote the perfect universal quotient on $B\times X$.
For the point $o\in X$ used to normalize the determinant, define
\[
 h=c_1^{\mathbb G_m}\bigl(\mathbf L i_o^*Q\bigr)
 \quad\text{in }\CH^1_{\mathbb G_m}(B),
 \qquad i_o:B\longrightarrow B\times X.
\]
The derived restriction in this definition is essential to make it an
operation on the perfect universal family, even where the quotient is
not locally free.  Let $q:B\times X\times X\to X\times X$ be projection,
and let the subscripts $(1),(2)$ mean pullback from the two copies of
$B\times X$.  Consider the nonlocalized equivariant class
\begin{equation}
 \Xi(t)=q_*\left(
 h^p\ch^{\mathbb G_m}_2(Q)_{(1)}
       \ch^{\mathbb G_m}_2(Q)_{(2)}
 \cap\bigl([B]^{\vir}\times[X\times X]\bigr)\right).
 \label{rel:equivariant-push}
\end{equation}
All insertions are Chern-class operations of perfect complexes, and $q$
is proper.  The total equivariant codimension of $\Xi(t)$ is $p+4-v$.
Since the torus acts trivially on $X\times X$,
\[
 \CH^*_{\mathbb G_m}(X\times X)
 =\CH^*(X\times X)[t].
\]
Consequently the coefficient of ordinary codimension two in the
localized class would have to be accompanied by
\begin{equation}
 t^{p+2-v}.
 \label{rel:parameter-power}
\end{equation}
Its exponent is negative, whereas~\eqref{rel:equivariant-push} has no
negative powers of $t$.  This component is therefore zero.  Virtual
localization~\cite{GP} expresses the same class as a sum of fixed
pushforwards divided by the moving Euler classes.  Thus their ordinary
codimension-two components sum to zero, and this remains true after
$t=1/2$.  Each individual fixed pushforward has the same total
equivariant codimension $p+4-v$: its ordinary codimension-two part
is therefore $t^{p+2-v}$ times a class independent of $t$.
Consequently changing $t$ to $-t$ multiplies each such contribution
by $(-1)^{p+2-v}$.  This termwise statement will allow us to compare
the two orders of a fixed decomposition in Section~\ref{sec:rankone}.

\subsection{A fixed contribution, including its signs}

We first write an individual fiber contribution without the multiplier
from the divisor virtual class.  Fix ordered lines
$D_1=T,D_2=H$ with $T+H=D$, and point lengths $n_i\geq0$ satisfying
\eqref{eq:length}.  Put
\[
 M=X^{[n_1]}\times X^{[n_2]},\qquad
 m_i=m_{D_i}=\frac{D_i(D_i-K)}2.
\]
The variables $x_i$ are the hyperplane classes on the divisor linear
systems; they can also be regarded as formal variables in the formula
below.  Write $\mathcal I_i$ for the universal ideals pulled to
$M\times X$, and $Z_i$ for their universal subschemes.  For a line $P$
let $P^{[n_i]}$ be its rank-$n_i$ tautological bundle.

For a vector bundle $E$ of rank $r$, our Euler convention is
\[
 \eu_u(E)=\sum_{j=0}^r c_j(E)u^{r-j}.
\]
For a virtual bundle it is extended multiplicatively whenever $u$ has
nonzero constant equivariant part.  Equivalently,
$\eu_u(E)=u^{\rk E}c_{1/u}(E)$, expanded in positive Chow degree.
The fixed point-obstruction factors are
\begin{equation}
 \mathcal E(x)=\prod_{i=1}^2
 \eu_{x_i}\left(\bigl((K-D_i)^{[n_i]}\bigr)^\vee\right).
 \label{rel:fixed-euler}
\end{equation}
Indeed $F_i=\bigl((K-D_i)^{[n_i]}\bigr)^\vee\otimes\mathcal O(1_i)$
on the indicated product or fiber, so~\eqref{rel:fixed-euler} is
$\prod_i\eu(F_i)$.

Define perfect $K$-classes on $M$ by
\[
 \begin{split}
 P_{ij}&=L_j-L_i+D_i,\\
 B_{ij}&=L_j-L_i+D_i-D_j,\\
 E_{ij}(P)&=[\RGamma(P)]-
 [\RHom(\mathcal I_i,\mathcal I_j(P))],
 \qquad i\ne j.
 \end{split}
\]
Here and below relative $\RHom$ includes derived pushforward along $X$.
The class $E_{ij}(P)$ is the standard Ext $K$-class of Carlsson--Okounkov
\cite[Section~1.2, equation~(2)]{CarlssonOkounkov}.  The moving ratio below
is derived here; their vertex-operator evaluation theorem is not used.
The inverse Euler class of the moving virtual normal bundle is
\begin{equation}
 \mathcal R(x)=\prod_{i\ne j}
 \frac{(w_j-w_i+x_i-x_j)^{\chi(B_{ij})}}
      {(w_j-w_i+x_i)^{\chi(P_{ij})}}
 \frac{\eu_{w_j-w_i+x_i}
   \left(\bigl((K-P_{ij})^{[n_i]}\bigr)^\vee\right)}
      {\eu_{w_j-w_i+x_i-x_j}\bigl(E_{ij}(B_{ij})\bigr)}.
 \label{eq:moving-ratio}
\end{equation}
We derive the signs explicitly.  On the fiber under consideration,
\[
 S'_i=L_i(-D_i)\mathcal I_i\otimes\mathcal O(-1_i).
\]
The ideal sequence and relative Serre duality give
\[
 [\RHom(\mathcal I_i,P)]
 =[\RGamma(P)]-
 \left[\bigl((K-P)^{[n_i]}\bigr)^\vee\right].
\]
The duality shift is $[-2]$, so it contributes with a positive sign
before subtraction by the ideal sequence.  Put $u=w_j-w_i$.  The moving
class for the ordered pair $(i,j)$, namely
$\RHom(S'_i,L_j)-\RHom(S'_i,S'_j)$, is therefore
\begin{align*}
 &[\RGamma(P_{ij})]_{u+x_i}
 -\left[\bigl((K-P_{ij})^{[n_i]}\bigr)^\vee\right]_{u+x_i}\\
 &\hspace{8mm}-[\RGamma(B_{ij})]_{u+x_i-x_j}
 +[E_{ij}(B_{ij})]_{u+x_i-x_j}.
\end{align*}
The subscript denotes the specified weight and line twist.  A constant
complex of vector spaces contributes the indicated weight raised to
its Euler characteristic.  Taking the inverse Euler class of this
four-term expression proves~\eqref{eq:moving-ratio}; in particular
$E_{ij}$ occurs in the denominator.  No individual cohomology dimension
of $P_{ij}$ or $B_{ij}$ remains in the result.

For the Chern-character calculation, temporarily read $w_1,w_2$ as the
degree-one equivariant classes $t,-t$: take the equivariant
codimension-two component before specializing $t=1/2$.
The restriction of the quotient Chern character to a marked surface
factor then follows from
\[
 \begin{aligned}
 \ch(Q_i)&=e^{L_i+w_i}
 \left(1-e^{-D_i-x_i}\ch(\mathcal I_i)\right),\\
 \ch(\mathcal I_i)
 &=1-[Z_i]+\text{terms of codimension at least three}.
 \end{aligned}
\]
Its first Chern class is $D_i+x_i$.  The determinant of the universal
point ideal is trivial, so derived restriction to $o$ adds no incidence
term to this first Chern class.  Line classes from $X$ restrict to zero
in $\CH^1(B)$, and the quotient has rank zero.  Thus $h$ restricts to
$x_1+x_2$.  In particular, fixing the determinant in the Picard scheme
does not make $h$ zero.  On this fixed product the determinant line of
the universal quotient is
$D\boxtimes\mathcal O(1_1+1_2)$, up to a constant line.
Its based normalization removes the factor from the parameter space,
whereas restriction of the unnormalized line to $o$ retains that factor.
This is exactly the distinction in Equation~\eqref{der:derived-fiber}.
The second Chern character on $M\times X$ is
\begin{equation}
 J=\sum_{i=1}^2
 \left([Z_i]+(L_i+w_i)(D_i+x_i)-\frac{(D_i+x_i)^2}{2}\right).
 \label{eq:insertion}
\end{equation}

Let $\pi:M\times X\times X\to X\times X$ be projection.  The
ordinary codimension-two fiber contribution is the following class,
where capping with $[M\times X\times X]$ in the pushforward is implicit:
\begin{equation}
 \begin{split}
 C_*(T,H;n_1,n_2)
 &=\Bigl[\,[x_1^{m_T}x_2^{m_H}]\\
 &\qquad\pi_*\bigl(\mathcal E(x)\mathcal R(x)(x_1+x_2)^p
 J_{(1)}J_{(2)}\bigr)\Bigr]_2.
 \end{split}
 \label{rel:fiber-contribution}
\end{equation}
The outer subscript means ordinary codimension two on $X\times X$.
All denominators are expanded at $x_1=x_2=0$ and in positive Chow
degree.  Their constant weights are $\pm1$, so these expansions are
well defined.  Only finitely many terms can contribute to the specified
$x$-coefficient and Chow degree.  A negative value of either $m_i$ means
that the contribution is zero.

For $p_g=q=0$, Proposition~\ref{prop:fixed-pg0} inserts
$x_i^{h^1(D_i)}\eu(F_i)$ on $|D_i|$.
Since $h^2(D_i)=0$ for effective $D_i$ and
$\dim|D_i|=m_i+h^1(D_i)$, projective integration is precisely the
coefficient extraction in~\eqref{rel:fiber-contribution}.
For the auxiliary setting, Proposition~\ref{prop:fixed-auxiliary}
first restricts to a fiber over the empty or canonical low divisor.
The low variable is $x_b=0$, its $m_b$ is zero, and the high projective
factor has dimension $m_a$.  The same formula applies, with multiplier
$1$ or $-s$, respectively.  This restriction precedes every calculation
with the insertions; there is no remaining integration over a Picard
parameter in~\eqref{rel:fiber-contribution}.

\subsection{Universality with both output factors retained}

Let $R_X\subset\CH^*(X)$ denote the graded subring generated by
$K,c_2(T_X),L_1,L_2,T,H$ and the identity.  We shall show that
\eqref{rel:fiber-contribution} lies in
\[
 \mathbb Q\cdot[\Delta_X]+\sum_{a+b=2}R_X^a\boxtimes R_X^b.
\]
We need an equality of Chow classes and a consistent numerical choice
of the coefficient of $\Delta_X$.  We therefore give the recursion
with the two output factors left unintegrated.

The reason for retaining these factors is already visible on $X^3$.
Let $p_{12}:X^3\to X^2$ retain the first two factors, let $p_i$ be
the individual projections, and write $\Delta_{ij}$ for the locus
where coordinates $i$ and $j$ agree.  For
$\gamma,\delta\in R_X$ homogeneous with total codimension two,
\begin{equation}\label{rel:three-slot-model}
 \begin{aligned}
 p_{12*}\bigl([\Delta_{13}]\cdot[\Delta_{23}]\bigr)
   &=[\Delta_X],\\
 p_{12*}\bigl([\Delta_{13}]\cdot p_3^*\gamma\cdot p_2^*\delta\bigr)
   &=\gamma\boxtimes\delta.
 \end{aligned}
\end{equation}
In the first identity both retained coordinates equal the third;
the small diagonal maps isomorphically to $\Delta_X$.
In the second, only the first coordinate is tied to the third;
the second remains independent.  These are Chow identities, before
any integration over the outputs.  The Hilbert-scheme recursion below
produces products of precisely such equality diagonals, together
with line and tangent Chern classes.  Their intersections will give
the same two alternatives, including the excess
factors when an equality is repeated.

\begin{lemma}[Consecutive Hilbert-scheme recursion]
\label{rel:hilbert-recursion}
Let $r>0$, and let $N_r=X^{[r-1,r]}$ be the scheme of nested
subschemes $Z'\subset Z$ of lengths $r-1,r$.  Write
\[
 \psi:N_r\longrightarrow X^{[r]},\quad
 \phi:N_r\longrightarrow X^{[r-1]},\quad
 \rho:N_r\longrightarrow X
\]
for the two forgetful maps and the residual-point map.  Then $N_r$ is
smooth, projective, irreducible of dimension $2r$;
$\psi_*[N_r]=r[X^{[r]}]$; and, with quotient-line convention,
\[
 \sigma=(\phi,\rho):N_r
 =\mathbb P(\mathcal I_{r-1})
 \longrightarrow X^{[r-1]}\times X.
\]
Put $\mathcal L=\mathcal O_{N_r}(1)$ and $\zeta=c_1(\mathcal L)$.
If $j:N_r\hookrightarrow N_r\times X$ is the graph of $\rho$, then
\begin{equation}
 0\longrightarrow(\psi\times\id_X)^*\mathcal I_r
 \longrightarrow(\phi\times\id_X)^*\mathcal I_{r-1}
 \longrightarrow j_*\mathcal L\longrightarrow0.
 \label{rel:universal-ideal-sequence}
\end{equation}
For a two-term vector bundle resolution
$0\to U_1\to U_0\to\mathcal I_{r-1}\to0$ on
$X^{[r-1]}\times X$, one has, for every $k\geq0$,
\begin{equation}
 \sigma_*(\zeta^k\cap[N_r])
 =\left[\frac{c(U_1^\vee)}{c(U_0^\vee)}\right]_k
 \cap[X^{[r-1]}\times X].
 \label{rel:projective-pushdown}
\end{equation}
All these identities remain valid after adjoining other Hilbert or
surface factors.
\end{lemma}

\begin{proof}
The geometric description and universal sequence are the consecutive
nested-scheme construction of~\cite[Section~1, equations~(3)--(6)]{EGL}.
A quotient $\mathcal I_{Z'}\twoheadrightarrow\mathbb C_x$ determines
the extra length-one subscheme and conversely.  In families its quotient
line is $\mathcal L$, giving the projectivization and
\eqref{rel:universal-ideal-sequence}.  The smoothness and dimension of
the consecutive nested scheme are the surface Hilbert-scheme facts
recorded there.  Over the open locus of $r$ distinct points, forgetting
which point is residual has degree $r$, proving the pushforward
formula for $\psi$.  The universal ideals are flat over their Hilbert
bases, so their ordinary pullbacks in this sequence agree with their
derived pullbacks.

Here is also the pushdown calculation, to fix its convention and sign.
The universal ideal has projective dimension at most one: it is flat
over the smooth Hilbert base and its surface fibers are torsion-free
ideals.  A sufficiently negative sum of ample line bundles surjects
onto it; the kernel is locally free, by the fiberwise projective
dimension bound and the local criterion of flatness.  This gives the
displayed global resolution.  Write $a=\rk U_1$, so $\rk U_0=a+1$,
and let $p:\mathbb P(U_0)\to X^{[r-1]}\times X$.
The inclusion $N_r\subset\mathbb P(U_0)$ is the zero locus of the
section of $p^*U_1^\vee\otimes\mathcal O(1)$ induced by $U_1\to U_0$.
Its codimension is $a$, because both $N_r$ and the base of $\sigma$
have dimension $2r$.  The ambient space is smooth, so the section is
regular and its top Chern class represents $[N_r]$.

For $\xi=c_1(\mathcal O_{\mathbb P(U_0)}(1))$ the quotient convention
gives
$p_*(\xi^{a+k})=[c(U_0^\vee)^{-1}]_k$ for $k\geq0$, and zero for
smaller exponents.  Hence the left side of
\eqref{rel:projective-pushdown} equals
\[
 \begin{split}
 p_*\bigl(\xi^k c_a(p^*U_1^\vee\otimes\mathcal O(1))\bigr)
 &=\sum_{b=0}^{\min(a,k)}c_b(U_1^\vee)
       [c(U_0^\vee)^{-1}]_{k-b}\\
 &=\left[\frac{c(U_1^\vee)}{c(U_0^\vee)}\right]_k.
 \end{split}
\]
This is also~\cite[Section~1, Lemma~1.1]{EGL}.  In particular, for
$r=2$ the map $\sigma$ is the blowup of $X\times X$ along its diagonal,
$\zeta=-E$ for its exceptional divisor, and the formula gives
$\sigma_*(\zeta^2)=-[\Delta_X]$.  Flat base change and the projection
formula give all the assertions with extra factors.
\end{proof}

\begin{lemma}[Reduction to decorated equality diagonals]
\label{rel:marked-reduction}
Fix a finite list of line bundles on $X$.  Consider a finite polynomial
of Chern characters of universal point ideals and their derived duals,
together with Chern classes of these lines and of $T_X$ pulled from
the surface slots and equality diagonals between those slots, on
\[
 X^{[n_1]}\times X^{[n_2]}\times X^\ell,
 \qquad \ell\geq2.
\]
Keep two specified surface factors as outputs and push forward over
the Hilbert factors and all other surface factors.  The resulting
Chow class is a finite rational linear combination of pushforwards of
products of equality diagonals and line and tangent classes on powers
of $X$.  The coefficients and the list of such products can be chosen
by one recursion independent of $X$ and of the chosen line bundles.
\end{lemma}

\begin{proof}
We induct on $n_1+n_2$.  When this sum is positive, choose a positive
length $r$ in a fixed order of the Hilbert factors and pull back along
$\psi$ of Lemma~\ref{rel:hilbert-recursion}, with all other factors
unchanged.  All spaces used here are smooth, so Chow pullbacks are
defined.  Projection formula and $\psi_*[N_r]=r[X^{[r]}]$ express the
original pushforward as $1/r$ times the pushforward of this pullback.

In every surface slot, replace the pulled-back ideal by
\eqref{rel:universal-ideal-sequence} in $K$-theory.  Its extra term is
the graph of the new residual point, with the line twist $\mathcal L$.
For that graph, Grothendieck--Riemann--Roch gives
\begin{equation}
 \ch(j_*\mathcal L)
 =j_*\bigl(e^{\zeta}\rho^*\td(T_X)^{-1}\bigr).
 \label{rel:graph-grr}
\end{equation}
Indeed the graph is a regular embedding of codimension two with
normal bundle $\rho^*T_X$.  It is the pullback of the diagonal under
$(\rho,\id_X):N_r\times X\to X\times X$.  This square is Tor
independent: the two local equations of the diagonal cut out the graph
regularly, in the same codimension.  Thus the graph term in
\eqref{rel:graph-grr} is a polynomial in $\zeta$, the equality diagonal
between the new point slot and the slot in use, and tangent classes.
For a perfect complex, the degree-$k$ Chern character of its derived
dual is $(-1)^k\ch_k$; derived duals therefore introduce no new kind
of class.

After this substitution the integrand is a polynomial in $\zeta$
whose coefficients pull back from
$X^{[r-1]}\times X$ and the unchanged factors.  Apply
\eqref{rel:projective-pushdown} to each power of $\zeta$.
Its right side is a universal polynomial in the Chern classes of
$\mathcal I_{r-1}$, equivalently in its Chern characters with rational
coefficients.  The Hilbert length has decreased by one; the added
surface factor records the residual point.  The two output factors
have not been pushed forward at any step.  Repeating the construction
ends with Hilbert schemes of length zero, whose ideals are the
structure sheaf.  Only surface factors, their diagonals, and the
specified decorations remain.  The choices of Hilbert-factor order,
slot order, and polynomial expansion order can be fixed once for all
inputs.  This proves the assertion.  We have adapted the geometric
recursion underlying the integral formula of~\cite[Section~3,
Proposition~3.1]{EGL}: here the two output factors remain unintegrated
throughout a sequence of Chow identities.
\end{proof}

The recursion has eliminated the Hilbert factors without integrating
the outputs. We now apply it to the localization integrand and use
the remaining diagonal intersections to select its coefficient of
\(\Delta_X\) by a rule independent of the surface.

\begin{proposition}[Chow universality of a fiber contribution]
\label{prop:universal}
The class in~\eqref{rel:fiber-contribution} admits a decomposition
\begin{equation}
 C_*(T,H;n_1,n_2)=c_*[\Delta_X]+\Gamma_*,
 \qquad
 \Gamma_*\in\sum_{a+b=2}R_X^a\boxtimes R_X^b.
 \label{rel:universal-decomposition}
\end{equation}
One can choose $c_*$ by a fixed universal rule.  Besides the integer
parameters and the chosen weights, this rule depends only on
\begin{equation}
 \int_X c_2(T_X),\qquad
 \int_X UV\quad
 (U,V\in\{K,L_1,L_2,T,H\}).
 \label{rel:numerical-data}
\end{equation}
It gives $c_*=0$ when $n_1=n_2=0$.  The same rule applies on every
smooth projective surface with the same data and $\chi(\mathcal O_X)=1$,
whether or not the lines $T,H$ are effective; effectivity is needed
only for realizing the formula as a fixed contribution.
\end{proposition}

\begin{proof}
We first put the integrand into the form required by
Lemma~\ref{rel:marked-reduction}.  For the surface projection
$\pi_X:M\times X\to M$, Grothendieck--Riemann--Roch gives
\begin{align}
 \ch\bigl(\RHom(\mathcal I_i,\mathcal I_j(P))\bigr)
 &=\pi_{X*}\bigl(\ch(\mathcal I_i^\vee)
          \ch(\mathcal I_j)e^P\td(T_X)\bigr),
 \label{rel:ext-grr}\\
 \ch(P^{[n_i]})
 &=\pi_{X*}\bigl((1-\ch(\mathcal I_i))e^P\td(T_X)\bigr).
 \label{rel:tautological-grr}
\end{align}
The dual in~\eqref{rel:ext-grr} is the derived dual.  These identities
hold in Chow, since the ideals are perfect and projection is smooth
and proper.  The constant complex $\RGamma(P)$ has $K$-class
$\chi(P)[\mathcal O_M]$, with
$\chi(P)=1+P(P-K)/2$.  The Chern classes in the Euler factors are
universal polynomials in the resulting Chern characters.  The
incidence insertion is $[Z_i]=-\ch_2(\mathcal I_i)$.
Thus, after the finite truncation required by
\eqref{rel:fiber-contribution}, every term is a polynomial in the
classes of the preceding Lemma and in pushforwards of such classes.

Products of surface pushforwards introduce only additional surface
slots.  More precisely, if $a,b$ are classes on $M\times X$, the
projection formula and flat base change identify
$\pi_{X*}(a)\pi_{X*}(b)$ with the pushforward of
$a_{(1)}b_{(2)}$ from $M\times X\times X$.  The same statement holds
over the base $M$ together with the two marked factors and can be
iterated for any finite product.  Apply
Lemma~\ref{rel:marked-reduction} to the resulting expressions.

It remains to perform the surface pushforwards.  This reduction will
both prove the Chow decomposition and specify the coefficient to be
compared across surfaces.  Represent a monomial in equality diagonals
by a graph $\mathcal G$: its vertices are the surface slots and each
factor $[\Delta_{uv}]$ is an edge between slots $u,v$, with repeated
factors recorded as repeated edges.  For each connected component $C$,
let $v_C,e_C$ count its vertices and edges.  Let
\[
 \iota:\prod_{C\text{ component of }\mathcal G}X
       \longrightarrow X^{\{\text{vertices of }\mathcal G\}}
\]
assign the same coordinate to all vertices in a component, and let
$\operatorname{pr}_C$ project its domain to the factor indexed by $C$.
Then
\begin{equation}\label{rel:graph-intersection}
 \prod_{uv\text{ edge of }\mathcal G}[\Delta_{uv}]
 =\iota_*\left(
    \prod_{C\text{ component of }\mathcal G}
    \bigl(\operatorname{pr}_C^*c_2(T_X)\bigr)^{e_C-v_C+1}
             \right).
\end{equation}
Indeed, a spanning tree in $C$ imposes $v_C-1$ transverse equalities,
giving its small diagonal.  Each of the remaining $e_C-v_C+1$ edges
imposes an equality already satisfied.  The diagonal
self-intersection formula contributes $c_2(T_X)$ for each such edge.
The exponent is zero for an isolated vertex as well as for a tree.
By projection formula, every line or tangent decoration is then
restricted to the copy of $X$ belonging to its component.

There are three possibilities for the subsequent pushforward:
\begin{enumerate}[label=\textup{(\roman*)}]
\item A component containing neither output is integrated over $X$.
It contributes zero unless its decoration has codimension two;
otherwise it contributes the intersection number of that decoration.
\item A component containing exactly one output maps isomorphically
to that output and leaves its decoration there.  When the outputs
belong to different components, their decorations therefore give an
external product of total codimension two, multiplied by the
intersection numbers from components of type~\textup{(i)}.
\item A component containing both outputs maps to their diagonal.
It contributes $\Delta_{X*}(\eta)$, multiplied by the same kind of
intersection numbers.  Since $\Delta_X$ already has codimension two,
the codimension-two part of the answer uses only
$\eta\in\CH^0(X)=\mathbb Q[X]$.  Thus this case gives a scalar
multiple of $[\Delta_X]$.
\end{enumerate}

Fix the recursion and edge-processing orders once for all inputs.
Carry out this reduction on the formal decorated graphs, replacing
components of type~\textup{(i)} by formal intersection numbers.
Define $c_*$ as the sum of the scalar coefficients in
case~\textup{(iii)}, and assign the terms with separated outputs to
$\Gamma_*$.  Only after this assignment do we evaluate the
intersection numbers and Chow classes on $X$.
Equation~\eqref{rel:graph-intersection} and the three cases prove
\eqref{rel:universal-decomposition} in Chow and specify the promised
choice of its coefficient.

Every codimension-two decoration integrated over a surface is a
linear combination of $c_2(T_X)$ and products of the lines in
\eqref{rel:numerical-data}, since $c_1(T_X)=-K$.
The scalar Euler exponents are determined by the same data through
Riemann--Roch.  Hence the rule gives the same $c_*$ on any two
surfaces with those numerical data.  We do not require the resulting
diagonal decomposition to be unique on an individual surface:
relations between its diagonal and external products are never
used to select $c_*$.

Finally, when both point lengths are zero, $M$ is a point,
$\mathcal E=1$, and $E_{ij}=0$.  There is no incidence term in $J$.
The two copies of $J$ involve only classes on their separate output
factors and the scalar variables $x_i,w_i$.  The graph procedure has
no edge joining the outputs, so it selects $c_*=0$.  None of this
formal calculation requires sections of $T$ or $H$, proving the last
assertion as well.
\end{proof}

All diagonal coefficients below use this fixed formal choice.
In particular, numerical comparison transports the coefficient of
the explicit fiber formula~\eqref{rel:fiber-contribution}; it does
not transport the geometry or multiplicity of the divisor component
from which that formula arose.

\subsection{The relation, cohomological tests, and its action}

The fixed contributions now have Chow decompositions with consistently
chosen diagonal coefficients. Summing them gives the required relation;
cohomological pairings will then test its coefficient, and its
correspondence action will control zero-cycles.

\begin{theorem}[Chow relation]
\label{thm:chow-relation}
For either surface setting, and any $p\geq0$ with $v>p+2$, sum the
classes~\eqref{rel:fiber-contribution} with the following multiplicities:
\begin{enumerate}[label=\textup{(\roman*)}]
 \item If $p_g=q=0$, sum over the effective ordered line pairs $T+H=D$
 and nonnegative point lengths satisfying~\eqref{eq:length}, with
 $m_T,m_H\geq0$.  Each fixed part has multiplier $1$.
 \item In the auxiliary setting, $D=3K$.  Sum over the ordered pairs
 $(3K,0),(0,3K)$ with multiplier $1$, and $(2K,K),(K,2K)$ with
 multiplier $-s$, with the admissible point lengths.  Here $s>0$ is
 the canonical-divisor coefficient of~\eqref{eq:canonical-cycle}.
\end{enumerate}
All sums are finite.  Writing $c$ for the same weighted sum of the
coefficients chosen in Proposition~\ref{prop:universal}, one obtains
\begin{equation}
 0=c[\Delta_X]+\Gamma
 \qquad\text{in }\CH^2(X\times X),
 \label{rel:chow-identity}
\end{equation}
where $\Gamma$ is a sum of external products of total codimension two.
\end{theorem}

\begin{proof}
The finite-type fixed locus has only finitely many discrete
decompositions and Hilbert polynomials, as in Section~\ref{sec:fixed}.
Proposition~\ref{prop:fixed-pg0} gives the first list and discards
negative $m_i$.  Proposition~\ref{prop:fixed-auxiliary} gives the second
list, its fiber restrictions, and its multiplicities.  The moving
factor and insertions were computed in~\eqref{eq:moving-ratio}
and~\eqref{eq:insertion}, so these are exactly the contributions to the
ordinary codimension-two part of localization of
\eqref{rel:equivariant-push}.  That part vanishes by
\eqref{rel:parameter-power}.  Substitute
\eqref{rel:universal-decomposition} for each contribution to obtain
\eqref{rel:chow-identity}.
\end{proof}

For the subsequent nonvanishing arguments, let
$\alpha\in H^2(X,\mathbb R)$ and write
$\alpha U=\int_X\alpha\smile U$ and
$\alpha^2=\int_X\alpha\smile\alpha$.
Recall $\pi_X:M\times X\to M$ and set
\[
 \mu_i(\alpha)=\pi_{X*}([Z_i]\smile\alpha),\qquad
 \mu_\alpha=\pi_{X*}(J\smile\alpha).
\]
Here $[Z_i]$ and $J$ denote their cycle classes; $\alpha$ need not be
algebraic.  Expansion of~\eqref{eq:insertion} gives the useful formula
\begin{equation}
 \mu_\alpha=\sum_{i=1}^2
 \left(\mu_i(\alpha)+(w_i-x_i)(\alpha D_i)
                         +x_i(\alpha L_i)\right).
 \label{rel:mu-formula}
\end{equation}
Integrating over the two output factors first and using projection
formula shows that
\begin{multline}
 \int_{X\times X}C_*(T,H;n_1,n_2)
                   \smile(\alpha\boxtimes\alpha)\\
 = [x_1^{m_T}x_2^{m_H}]\int_M
 \mathcal E(x)\mathcal R(x)(x_1+x_2)^p\mu_\alpha^2.
 \label{rel:pairing}
\end{multline}
The integral takes the top ordinary cohomological degree on $M$.
It automatically selects precisely the ordinary codimension-two
part on the left.  In particular, no algebraicity assumption on
$\alpha$ is used.

If $\alpha$ is orthogonal to $K,L_1,L_2,T,H$, then
\begin{equation}
 \int_{X\times X}C_*(T,H;n_1,n_2)
                  \smile(\alpha\boxtimes\alpha)
 =c_*\alpha^2.
 \label{rel:orthogonal-test}
\end{equation}
Indeed, an external product of codimensions $(0,2)$ or $(2,0)$ pairs
to zero by degree, and one of codimensions $(1,1)$ pairs to the
product of two intersections with $\alpha$.  Those intersections
vanish because the codimension-one part of $R_X$ is spanned by the
listed line classes.  The diagonal pairs to $\alpha^2$.
This is a test for the chosen coefficient of a relation already
proved in Chow.

\begin{corollary}
\label{rel:correspondence}
If $p_g=q=0$ and the coefficient $c$ in
\eqref{rel:chow-identity} is nonzero, then
$\CH_0(X)^0\otimes\mathbb Q=0$.
\end{corollary}

\begin{proof}
For $Z\in\CH^2(X\times X)$, use its correspondence action
\[
 Z_*(z)=\operatorname{pr}_{2*}
       \bigl(\operatorname{pr}_1^*z\cdot Z\bigr).
\]
The diagonal acts as the identity.  If $Z=a\boxtimes b$ has total
codimension two and $a$ has positive codimension, then
$z\cdot a=0$ on the surface for every zero-cycle $z$, by dimension.
If $a$ has codimension zero, connectedness of $X$ gives
$a=\lambda[X]$ and $Z_*(z)=\lambda\deg(z)b$.
Every summand of $\Gamma$ therefore annihilates degree-zero
zero-cycles.  Acting by~\eqref{rel:chow-identity} gives $cz=0$ in
$\CH_0(X)\otimes\mathbb Q$.  Since $c\ne0$, this implies $z=0$.
\end{proof}
